% ==============================================================================
% Plantilla Institucional de Trabajo Terminal — UPIIZ IPN
% Archivo: capitulos/protocolo/06-descripcion-trabajo-propuesto.tex
% ==============================================================================
% ESTRUCTURA NORMATIVA DE PROTOCOLO: CAPÍTULO 6 - DESCRIPCIÓN DEL TRABAJO PROPUESTO
% ==============================================================================
% LÍMITE INSTITUCIONAL: Máximo 6 páginas.
%
% Requisitos normativos (Anexo 1):
% 1. Explicación general del sistema a desarrollar.
% 2. Alcance.
% 3. Delimitaciones.
% 4. Diagramas IDEF-0 A-0 y A0.
% 5. Diagramas de flujo.
% 6. Actividades necesarias para cumplir los objetivos.
% 7. Áreas de conocimiento involucradas de acuerdo con la Línea de Trabajo.
%
% NOTA SOBRE DIAGRAMAS:
% Los diagramas IDEF-0 y diagramas de flujo pueden generarse con cualquier herramienta
% externa e insertarse como figuras mediante \includegraphics. No se impone ninguna
% herramienta o paquete específico.
% ==============================================================================

\chapter{Descripción del trabajo propuesto}
\label{cap:proto-descripcion-propuesta}

% [INSTRUCCIÓN PARA EL ALUMNO]:
% Desarrolle este capítulo en un máximo de 6 páginas cubriendo todos los puntos normativos:

\section{Explicación general del sistema a desarrollar}

[Describa de forma general y comprensible el sistema, máquina, proceso, dispositivo o software que se proyecta desarrollar, explicando su funcionamiento conceptual y propósito global].

\section{Alcance}
\label{sec:proto-alcance}

[Establezca con precisión el alcance formal del proyecto: qué funciones desempeñará el sistema, qué componentes se entregarán y hasta qué etapa de desarrollo o validación se llegará].

\section{Delimitaciones}
\label{sec:proto-delimitaciones}

[Puntualice las fronteras y restricciones del proyecto: qué aspectos quedan fuera del desarrollo, condiciones operativas que no se contemplan y límites de aplicación].

\section{Diagrama funcional IDEF-0 Nivel A-0}
\label{sec:proto-idef0-a-0}

% [INSTRUCCIÓN PARA EL ALUMNO]:
% Presente el diagrama de contexto IDEF-0 Nivel A-0 (función principal del sistema
% con sus entradas, controles, salidas y mecanismos):

\begin{figure}[htbp]
    \centering
    % Inserte aquí su figura generada externamente:
    % \includegraphics[width=0.85\textwidth]{figuras/proyecto/idef0-a-0.pdf}
    \fbox{\parbox[c][5cm][c]{0.85\textwidth}{\centering\textbf{[Espacio para Diagrama IDEF-0 Nivel A-0 (Contexto)]}\\[0.5em]
    \small Inserte aquí la figura generada externamente mediante:\\
    \texttt{\textbackslash includegraphics[width=0.85\textbackslash textwidth]\{figuras/proyecto/idef0-a-0.pdf\}}}}
    \caption{Diagrama funcional de contexto IDEF-0 Nivel A-0 del sistema a desarrollar.}
    \label{fig:proto-idef0-a-0}
\end{figure}

\section{Diagrama funcional IDEF-0 Nivel A0}
\label{sec:proto-idef0-a0}

% [INSTRUCCIÓN PARA EL ALUMNO]:
% Presente el diagrama de descomposición funcional IDEF-0 Nivel A0 (actividades
% principales del sistema y sus interconexiones):

\begin{figure}[H]
    \centering
    % Inserte aquí su figura generada externamente:
    % \includegraphics[width=0.90\textwidth]{figuras/proyecto/idef0-a0.pdf}
    \fbox{\parbox[c][6cm][c]{0.85\textwidth}{\centering\textbf{[Espacio para Diagrama IDEF-0 Nivel A0 (Descomposición Funcional)]}\\[0.5em]
    \small Inserte aquí la figura generada externamente mediante:\\
    \texttt{\textbackslash includegraphics[width=0.9\textbackslash textwidth]\{figuras/proyecto/idef0-a0.pdf\}}}}
    \caption{Diagrama de descomposición funcional IDEF-0 Nivel A0 del sistema.}
    \label{fig:proto-idef0-a0}
\end{figure}

\section{Diagramas de flujo}
\label{sec:proto-diagramas-flujo}

% [INSTRUCCIÓN PARA EL ALUMNO]:
% Inserte el diagrama de flujo correspondiente al proceso, algoritmo o secuencia de operación:

\begin{figure}[htbp]
    \centering
    % Inserte aquí su figura generada externamente:
    % \includegraphics[width=0.75\textwidth]{figuras/proyecto/diagrama-flujo.pdf}
    \fbox{\parbox[c][5.5cm][c]{0.8\textwidth}{\centering\textbf{[Espacio para Diagrama de Flujo]}\\[0.5em]
    \small Inserte aquí el diagrama de flujo mediante:\\[0.3em]
    \texttt{\textbackslash includegraphics[width=0.75\textbackslash textwidth]}\\[0.1em]
    \texttt{\{figuras/proyecto/diagrama-flujo.pdf\}}}}
    \caption{Diagrama de flujo del proceso o secuencia de operación prevista.}
    \label{fig:proto-diagrama-flujo}
\end{figure}

\section{Actividades necesarias para cumplir los objetivos}
\label{sec:proto-actividades}

[Desglose la lista secuencial de actividades técnicas indispensables que el equipo de trabajo deberá llevar a cabo para dar cumplimiento a los objetivos planteados].

\section{Áreas de conocimiento involucradas de acuerdo con la Línea de Trabajo}
\label{sec:proto-areas-conocimiento}

% Conforme al Anexo 1, todas las Líneas de Trabajo deben contemplar las 4 áreas
% formativas de la Mecatrónica, con el grado de énfasis correspondiente a la línea seleccionada:
% Línea de Trabajo: \NumeroLineaTrabajo\ - \NombreLineaTrabajo

[Detalle la participación de las áreas de conocimiento involucradas en el proyecto, de acuerdo con la Línea de Trabajo seleccionada (\NumeroLineaTrabajo) y las cuatro áreas formativas de la disciplina: mecánica, electrónica, control y programación].
